% ch00-profiling.tex — 性能分析先于优化：nsys 与 ncu（RFC-C0）
% 工具方法论章：无对应源码区间，无编译宏限定，全 arch 通用
% 本章无 device kernel：book/tests/ 不设 ch00 测试；实战数据引自作者参与的 CUDA kernel 开发 NCU/NSYS 分析记录
\chapter{性能分析先于优化：nsys 与 ncu}\label{ch:00}

\section{本章导读}

在动手写第一个 kernel 之前，本书想先立一条工程信条：\textbf{profile 永远都是
最重要的}。优化 GPU kernel 的正确循环不是「改了跑一下看快没快」，而是
「先测量、再假设、改一处、再测量」。凭直觉的优化是赌博：赢了是运气，输了
连原因都找不到。这一章专门讲两件测量工具——\textbf{Nsight Systems}（nsys，
全局时间线）与 \textbf{Nsight Compute}（ncu，单 kernel 计数器）——的原理、
命令行用法与核心指标的含义，并给出一张「症状 $\to$ 证据 $\to$ 结论 $\to$
对策」的瓶颈判断决策树。

本章不含 kernel（与 ch01 及 Part IV 的原理章 ch20--22 一样，是纯理论/工具章）。没有新代码可讲，但有大量
\textbf{实战教训}：它们全部来自真实的 CUDA kernel 开发记录——寄存器压力、spill、warp stall、Tensor Core 饱和度、
L1/L2 命中率、bank conflict，每一个指标都配一个真实的「指标驱动优化」或
「指标否决优化」案例。后续所有章节的「性能分析」节都默认读者已掌握本章的
工具与词汇表。

\textbf{前置}：无（与 ch01 互为入口，roofline 词汇见 ch01）。
\textbf{源码区间}：无（工具方法论章）。\textbf{编译宏与架构限定}：无。

\section{问题与动机：凭感觉的优化是赌博}

先看三个真实的开发反例，说明「不 profile 就优化」会栽在哪里。

\textbf{反例一：感觉是装载瓶颈，其实是依赖链}。一个 sm\_89 上的 fp8
attention kernel（persist-D 结构）对标业界最快实现慢 28\%（909.4$\mu$s vs
710$\mu$s）。直觉方案是「加深装载流水、砍指令、缩 tile」，全部实测无效；
NCU 的 stall 剖析结果一锤定音：\texttt{math\_\allowbreak{}pipe\_\allowbreak{}throttle} 1.32
主导、\texttt{wait} 0.82、\texttt{long\_\allowbreak{}scoreboard} 仅
0.04——\textbf{装载根本不是瓶颈}，瓶颈是单 CTA 内 softmax$\to$PV 的串行
依赖链，唯一有效的改动是拆链（2 CTA/SM），实测 $-5\%$。

\textbf{反例二：感觉主 kernel 慢，其实差距在别处}。e2e 落后时先盯着主 kernel
跑 ncu，做完整套指标分析才发现主 kernel 已占 e2e 时间 86--92\%、前处理链
已优于对照——一次 nsys 的时间线本可以三分钟得出同样结论，避免在错误的
kernel 上浪费半天 ncu 分析。

\textbf{反例三：感觉融合更快，其实更慢}。把 softmax 的 bias 注入从「分离
kernel」改成「单 pass 融合 + smem tile」，直觉上少一次读写应该更快，实测反而
慢 4.2ms：NCU 显示融合把 LDS（shared load）拉进了 consumer 关键路径，
\texttt{sleeping} stall 从 3.5\% 暴涨到 22.3\%。指标否决了这次「优化」。

三个反例的共同教训：GPU 上的性能瓶颈极少长在直觉指向的地方。成熟的优化
工作流会把它固化成一条硬规则：\textbf{没有 NCU 数据，不改 kernel}；
每条优化结论必须引用具体 metric 数值作为证据。本章就把这套「证据
从哪里来、怎么读」讲透。

\section{工具链总览与顺序方法论}\label{sec:ch00-tools}

\subsection{两个工具回答两个不同的问题}

nsys 与 ncu 的采集原理完全不同：

\begin{itemize}
\item \textbf{nsys（Nsight Systems）}：\textbf{trace} 型。在驱动层记录 API
  调用、kernel launch、内存拷贝、流与事件的时间戳，拼出一张全局时间线。
  它回答「\textbf{时间花在哪}」：哪个 kernel 占比最大、kernel 之间有没有
  空隙（gap）、CPU launch 是否跟不上、多流是否真的在重叠。代价低——不需要
  GPU 性能计数器权限，任何环境都能跑，被测程序几乎不失真。
\item \textbf{ncu（Nsight Compute）}：\textbf{counter replay} 型。要读一个
  kernel 的成百上千个硬件计数器（SM 吞吐、cache 命中、warp stall……），
  而一遍执行采不完这么多计数器，于是 ncu 把 kernel \textbf{回放多遍}，
  每遍采一组；为了数据可比还会锁定时钟。它回答「\textbf{这个 kernel 内部
  慢在哪}」。代价高——单 kernel 深挖、时间被串行化与放大，\textbf{ncu 报告
  里的 kernel 耗时不能当 bench 时间用}。
\end{itemize}

一句话分工：nsys 是体检报告（哪里有病），ncu 是专科检查（病灶细胞级切片）。

\subsection{顺序：先 nsys，后 ncu}

顺序不是习惯问题，是成本问题：

\begin{enumerate}
\item \textbf{nsys 先行}：一次采集拿到全部 kernel 的时间占比排序。如果主
  kernel 占 86\%，那么辅助 kernel 再难看也只有 14\% 的优化上限——先确认
  「该优化谁」，再谈「怎么优化」。nsys 还能立刻暴露三类结构性问题：launch
  gap（CPU 跟不上，一个真实案例 wall-GPU 差 129$\mu$s vs 对照 50$\mu$s，
  病根在 dispatch 而非 kernel）、该并行的流没重叠、意外的 fallback kernel。
\item \textbf{ncu 跟进}：只对 nsys 排出的 Top-1/Top-2 kernel 做
  \texttt{-{}-set full} 深挖，读 roofline、occupancy、stall、memory 各维
  指标，形成假设。
\item \textbf{单变量修改 $\to$ 复测}：每次只改一处，re-profile 后用
  \texttt{-{}-diff} 对照前后指标；无效或负优化立即回退。
\end{enumerate}

反着来的代价见本章开头的反例二：先 ncu 后 nsys，等于对着一个可能不重要的
kernel 做了半天细胞级切片。

\begin{figure}[H]
\centering
\begin{tikzpicture}[
  font=\footnotesize,
  box/.style={draw, rounded corners=2pt, align=center, inner sep=3pt,
    font=\scriptsize, text width=2.2cm, minimum height=1.0cm},
  ar/.style={-{Stealth[length=1.8mm]}, semithick}
]
  \node[box, text width=1.3cm] (e2e) at (0.95, 0) {e2e\\落后};
  \node[box, fill=gray!15] (nsys) at (3.4, 0)
    {nsys 全局时间线\\\texttt{-t cuda} 单域\\kern\_sum 占比排序};
  \node[box] (top) at (6.35, 0)
    {Top-1/2 kernel\\主 kernel 占 86--92\%\\$\to$ 辅助链上限仅 14\%};
  \node[box, fill=gray!40] (ncu) at (9.3, 0)
    {ncu 单 kernel 深挖\\\texttt{-{}-set full}\\\texttt{regex:} 名字前缀};
  \node[box] (judge) at (12.35, 0)
    {指标判读\\（接决策树\\图~\ref{fig:0-1}）};
  \draw[ar] (e2e.east) -- (nsys.west);
  \draw[ar] (nsys.east) -- (top.west);
  \draw[ar] (top.east) -- (ncu.west);
  \draw[ar] (ncu.east) -- (judge.west);
  \node[box, fill=gray!65] (modi) at (12.35, -2.35)
    {单变量修改\\每次只改一处};
  \node[box, fill=gray!65] (diff) at (8.9, -2.35)
    {复测 \texttt{-{}-diff}\\前后指标对照};
  \node[box] (keep) at (5.2, -2.35) {有效 $\to$ 保留};
  \draw[ar] (judge.south) -- (modi.north);
  \draw[ar] (modi.west) -- (diff.east);
  \draw[ar] (diff.west) -- (keep.east)
    node[midway, above, font=\scriptsize] {有效};
  \node[box, dashed] (revert) at (8.9, -4.1)
    {无效 / 负优化\\$\to$ 回退};
  \draw[ar] (diff.south) -- (revert.north)
    node[midway, right, font=\scriptsize] {无效};
  \draw[ar, dashed] (revert.east) -- (14.1, -4.1) -- (14.1, 0) -- (judge.east);
  \node[font=\scriptsize, anchor=west, text=black!60] at (14.25, -2.0)
    {回到判读};
\end{tikzpicture}
\caption{nsys $\to$ ncu $\to$ 单变量修改 $\to$ 复测 \texttt{-{}-diff} 的闭环工作流：nsys 先给全局占比与优化预算（主 kernel 占 86--92\%，辅助链上限仅 14\%，对应 \S\ref{sec:ch00-cases} 例一）；ncu 只深挖 Top-1/2 kernel；每次只改一处，复测后无效立即回退、回到指标判读。泳道灰度：nsys 15\%、ncu 40\%、修改/复测 65\%。}
\label{fig:0-2}
\end{figure}

\section{nsys 命令行实战}\label{sec:ch00-nsys}

\subsection{最小采集}

原则：\textbf{只采需要的域、只采最短的窗口}。对「看 kernel 占比」这个目标，
一条 trace 域 \texttt{-t cuda} 就够，CPU 采样与上下文切换全部关掉，warmup
与 repeat 压到最小：

\begin{lstlisting}[style=console, title={nsys 最小采集（kernel 时间占比场景）}]
nsys profile -t cuda --sample=none --cpuctxsw=none \
  -o .tmp/ch00/report.nsys-rep python bench.py
\end{lstlisting}

\subsection{stats 报表：两个最常用的 report}

采集产物是 \texttt{.nsys-rep}，用 \texttt{nsys stats} 导出统计：

\begin{lstlisting}[style=console, title={nsys stats 两个高频 report}]
nsys stats -r cuda_gpu_kern_sum report.nsys-rep   # per-kernel 时间占比（最常用）
nsys stats -r cuda_gpu_trace   report.nsys-rep   # per-launch 时间线（看 gap/流重叠）
\end{lstlisting}

\texttt{cuda\_\allowbreak{}gpu\_\allowbreak{}kern\_\allowbreak{}sum} 按 kernel 聚合总时长与占比，直接给出「该优化
谁」的答案；\texttt{cuda\_\allowbreak{}gpu\_\allowbreak{}trace} 逐 launch 列时间戳，用来看相邻 kernel
间的空隙与多流 overlap。首次对旧 report 运行时报
\texttt{Existing SQLite export found} 时，加 \texttt{-{}-force-export=\allowbreak{}true}
强制重新导出。

\subsection{体积控制：一次 500GB 事故}

nsys 的中间数据默认写到 \texttt{/\allowbreak{}tmp/\allowbreak{}nvidia/\allowbreak{}nsight\_\allowbreak{}systems}，体积 $\approx$
单事件字节数 $\times$ 事件数，而事件数随时长线性增长。对大模型推理（多模型链、大 seqlen、多次 repeat）无脑
全量采集，曾把 \texttt{/tmp} 写到 \textbf{500GiB+} 直到磁盘写满系统卡死
（真实事故，2026-09）。
三条铁律：

\begin{enumerate}
\item \textbf{输出与临时文件落项目目录}：显式设 \texttt{TMPDIR} 指向项目
  \texttt{.tmp/} 下的目录并用 \texttt{-o} 指定输出，事后清理。注意 TMPDIR
  路径要短——过长的路径会撑爆 Unix domain socket 的 107 字节上限（Ray 的
  plasma store 直接构造失败）。
\item \textbf{禁止对大模型 e2e 无脑全量采}：repeat 降到 0--1，单域
  \texttt{-t cuda}，关采样。
\item \textbf{优先 torch.profiler}：只想看 kernel 时间表时，
  \texttt{torch.\allowbreak{}profiler} 输出 chrome trace 落在 \texttt{.tmp/}，零外部临时
  文件风险；nsys 只在必须看时间线/同步事件时出场。
\end{enumerate}

\subsection{窗口化采集：cudaProfilerApi}

只想采推理中的一小段（比如 50 个 denoise step 中的 1 步）时，用
\texttt{-{}-capture-range=\allowbreak{}cuda\allowbreak{}Profiler\allowbreak{}Api} 配合代码里的开关，把采集窗口钉在
目标区间：

\begin{lstlisting}[style=console, title={窗口化采集（Python 侧）}]
nsys profile -t cuda --capture-range=cudaProfilerApi \
  --sample=none --cpuctxsw=none -o .tmp/ch00/win python bench.py
# bench.py 内：
#   torch.cuda.profiler.start()   # 窗口开始
#   ... 目标 step ...
#   torch.cuda.profiler.stop()    # 窗口结束
\end{lstlisting}

\section{ncu 命令行实战}\label{sec:ch00-ncu}

\subsection{基础采集与 section}

\begin{lstlisting}[style=console, title={ncu 全量采集与 CSV 导出}]
ncu --set full -o .tmp/ch00/report_v1 --target-processes all ./kernel_app
ncu --import .tmp/ch00/report_v1.ncu-rep --page raw --csv > report_v1.csv
\end{lstlisting}

\texttt{-{}-set full} 采全部 section（roofline、occupancy、warp state、
memory workload、instruction mix……），单 kernel 回放多遍、耗时可观；快速
确认可用默认 basic set。交互式 \texttt{ncu-ui} 看报告，命令行用
\texttt{-{}-page details}（rule-based 汇总）或 \texttt{-{}-page raw}（原始
metric，见下）。

\subsection{-{}-kernel-name 必须用 regex: 前缀}

C++ 模板 kernel 的完整名带一长串后缀（如
\texttt{persist\_\allowbreak{}d\_\allowbreak{}ws\_\allowbreak{}fwd\_\allowbreak{}cute\_\allowbreak{}fp8\_\allowbreak{}sm120}，实际实例化名比这更长）。
\texttt{-{}-kernel-name} 默认\textbf{精确匹配}，拿「名字的前缀」去匹配会报
\texttt{No kernels were profiled} 且生成空 report。必须加 \texttt{regex:}
前缀走部分匹配：

\begin{lstlisting}[style=console, title={kernel 名过滤：regex 前缀}]
ncu --set full -o rep --kernel-name regex:persist_d_ws_fwd_cute_fp8 \
  --target-processes all ./app     # OK：部分匹配
ncu --kernel-name persist_d_ws_fwd_cute_fp8_sm120 ...  # 错：精确匹配，0 命中
\end{lstlisting}

\subsection{-{}-page raw 提取原始 metric}

ncu 的 \texttt{-{}-csv} 默认输出是 \textbf{rule-based} 的 section 汇总（按
诊断规则组织），不是逐个原始计数器。要拿具体 stall reason 的数值，必须走
\texttt{-{}-page raw}。raw 输出里每个 metric 是「名字行 + 单位行 + 数值行」
的多行结构，用 \texttt{grep -A3} 抓：

\begin{lstlisting}[style=console, title={从 raw page 提取 stall 指标}]
ncu --import rep.ncu-rep --page raw 2>&1 | \
  grep -A3 "warps_issue_stalled_\(wait\|math_pipe_throttle\|sleeping\)"
\end{lstlisting}

\subsection{权限限制与同架构代理 GPU}

读性能计数器需要驱动授权。容器/云主机上常见
\texttt{/proc/\allowbreak{}driver/\allowbreak{}nvidia/\allowbreak{}params} 里
\texttt{Rm\allowbreak{}Profiling\allowbreak{}Admin\allowbreak{}Only} 标志为 1，
此时 \texttt{-{}-set full} 直接报 \texttt{ERR\_\allowbreak{}NVGPUCTRPERM}，且容器内无法
reload 驱动。可行解法是\textbf{同架构代理}：找一块相同微架构的本地 GPU 跑
ncu——例如用本地 RTX PRO 5000（sm\_120a）代理远端 RTX 5090（sm\_120）。
stall reasons、指令 mix、occupancy 这类\textbf{架构级}结论可以迁移，绝对
性能数字（SM 数、频率、L2 容量）不可迁移；ncu 验证完架构结论后，回目标
硬件跑 timing bench 收数。

\subsection{replay 的干扰与 A/B 纪律}

ncu 锁频 + 多遍回放 + 串行执行，因此\textbf{报告内耗时只用于相对比较}。
做优化 A/B 时回到真实 bench，并遵守两条纪律：

\begin{itemize}
\item \textbf{校验公共参照}：许多 bench 会同时输出被测实现与公共参照
  （如 PyTorch SDPA）两组时间；对比前先核对两轮参照是否吻合（$\pm$1\%）。
  被 GPU 级干扰（并发进程/时钟态）污染的一轮，两者会\textbf{同比例膨胀}
  （实测约 $2\times$），直接对比会得出「$-50\%$ 全面加速」这种物理上不可能
  的假结论。
\item \textbf{min-of-N 口径}：单轮抖动大的实现（如 tma 变体）取 N 次最小值
  对比；归因 race 类问题必须 stash $\to$ 重建 $\to$ 同 probe 的 HEAD 对照，
  单次「好结果」不可信。
\end{itemize}

\section{核心指标字典}\label{sec:ch00-metrics}

以下每个指标按「含义 / 怎么读 / 异常意味着什么 / 对策」展开。数值口径以
sm\_120 上的实测为参照。

\subsection{SpeedOfLight：SM\% 与 MEM\%（roofline 定位）}

\texttt{sm\_\allowbreak{}\_\allowbreak{}throughput}（SM\%）与 \texttt{gpu\_\allowbreak{}\_\allowbreak{}compute\_\allowbreak{}memory\_\allowbreak{}throughput}
（MEM\%）分别是计算管线与内存管线的峰值利用率，是 NCU 报告的第一行。
两者把 kernel 放上 ch01 的 roofline：SM\% 高 MEM\% 低是 compute-bound，
反之 memory-bound；\textbf{双低}（如 SM/MEM 都在 30\% 附近）则是
latency-bound——管线都没喂满，问题在并行度或依赖链，这是 roofline 图上
看不出的第三态，必须回到 stall 指标找证据。

\subsection{Occupancy：理论、达到与三类限制源}

区分两个数：\textbf{理论 occupancy}（kernel 资源能驻留多少 warp，受三类
限制源钳制——寄存器、smem、block/warp 上限）与\textbf{达到 occupancy}
（\texttt{sm\_\allowbreak{}\_\allowbreak{}warps\_\allowbreak{}active}，实际运行均值，受 tail effect 与负载不均
拖累）。NCU 的 Occupancy section 会明确指出 \texttt{launch\_\allowbreak{}\_\allowbreak{}occupancy\_\allowbreak{}limit\_\allowbreak{}registers} /
\texttt{\_\allowbreak{}\_\allowbreak{}shared\_\allowbreak{}mem} / \texttt{\_\_blocks} 哪个是当前限制源。一个
fp8 attention kernel 的实例：REG:255 直接把 \texttt{Block Limit Registers} 压到 1——
1 CTA/SM、12 warps（4 producer + 8 consumer，共 384 线程），理论 occupancy
25\%（12/48），此时想提并行度只有降寄存器或换结构
（见 \S\ref{sec:ch00-cases} 例二）。注意 ch01 的告诫：occupancy 是手段不是
KPI，GEMM 类 kernel 用大 tile + 低 occupancy + 高 ILP 换性能是常态。

\begin{figure}[H]
\centering
\begin{tikzpicture}[font=\footnotesize,
  lab/.style={font=\scriptsize, inner sep=1pt}]
  \draw[-{Stealth[length=1.8mm]}] (0, 0) -- (13.3, 0);
  \draw[-{Stealth[length=1.8mm]}] (0, 0) -- (0, 6.75);
  \foreach \x/\t in {1.44/32, 2.88/64, 5.76/128, 8.64/192, 11.52/256}
    {
      \draw (\x, 0) -- (\x, -0.08);
      \node[lab, anchor=north] at (\x, -0.1) {\t};
    }
  \foreach \y/\t in {0.92/8, 1.84/16, 2.76/24, 3.68/32, 4.60/40}
    {
      \draw (0, \y) -- (-0.08, \y);
      \node[lab, anchor=east] at (-0.1, \y) {\t};
    }
  \draw (0, 5.52) -- (-0.08, 5.52);
  \node[lab, anchor=east, font=\scriptsize\bfseries] at (-0.1, 5.52) {48};
  \draw[black!30, semithick] plot[smooth] coordinates
    {(1.80,5.89) (2.16,4.91) (2.52,4.21) (2.88,3.68) (3.24,3.27) (3.60,2.94)
     (4.32,2.45) (5.04,2.10) (5.76,1.84) (7.20,1.47) (8.64,1.23) (10.08,1.05)
     (11.52,0.92)};
  \draw[black!45, semithick] (0.3, 3.68) -- (7.9, 3.68);
  \draw[black!60, semithick] (0.3, 5.52) -- (11.8, 5.52);
  \draw[very thick]
    (0.3,3.68) -- (2.88,3.68) -- (2.88,2.76) -- (3.83,2.76) -- (3.83,1.84)
    -- (5.76,1.84) -- (5.76,0.92) -- (11.52,0.92);
  \node[lab, anchor=south west] at (4.7, 5.72) {block/warp 槽位上限：48 warps};
  \node[lab, anchor=south west, align=left] at (3.3, 3.9)
    {smem 约束（示例：\\smem/block=32KB，T=256）\\$\to$ 上限 32 warps};
  \node[lab, anchor=west, align=left] at (11.68, 0.92)
    {寄存器约束\\$65536/(32r)$};
  \node[lab, anchor=west, align=left] at (11.68, 2.35)
    {粗黑阶梯 = min\\（整 block 量化，\\T=256 $\to$ 8 warps/级）};
  \fill (11.48, 1.38) circle (0.07);
  \draw[black!60, thin] (9.8, 4.1) -- (11.42, 1.47);
  \node[lab, anchor=north west, align=left] at (8.45, 5.25)
    {实测点（WS 结构，例二）：\\consumer REG 255 / producer 32\\
     1 CTA/SM = 12 warps（384T）\\理论 occupancy 25\%（12/48）};
  \node[lab, anchor=north] at (6.65, -0.42) {每线程寄存器数 $r$（regs/thread）};
  \node[lab, rotate=90, anchor=south] at (-0.95, 3.3) {可驻留 warp/SM};
\end{tikzpicture}
\caption{理论 occupancy 的三限制源取 min（sm\_120，48 warp 槽位口径）：寄存器约束 $65536/(32r)$（warp 粒度简化）随每线程寄存器数 $r$ 下降；smem 与槽位为水平约束（示例：smem/block=32KB、T=256 $\to$ 32 warps）。粗黑阶梯为按整 block（T=256，8 warps/block）量化的 min 包络；实测点为 WS 结构（producer 32 / consumer 232 regs）1 CTA/SM = 12 warps（384T），理论 occupancy 25\%（12/48）。}
\label{fig:0-3}
\end{figure}

\subsection{寄存器与 spill：255 墙}

\texttt{launch\_\allowbreak{}\_\allowbreak{}registers\_\allowbreak{}per\_\allowbreak{}thread} 给出每线程寄存器数，硬上限
\textbf{255}（每线程）。超限时编译器把变量溢出到 local memory（SASS 里表现为
\texttt{LDL}/\texttt{STL} 指令，NCU 里看
\texttt{l1tex\_\allowbreak{}\_\allowbreak{}t\_\allowbreak{}bytes\_\allowbreak{}pipe\_\allowbreak{}lsu\_\allowbreak{}mem\_\allowbreak{}\allowbreak{}local\_\allowbreak{}op\_\allowbreak{}\_\allowbreak{}ld/\allowbreak{}st.sum} 是否非零）。
编译期排查用 \texttt{-{}-ptxas-options=\allowbreak{}-v} 与
\texttt{cuobjdump -{}-dump-resource-usage} 的 STACK 字节。两个实战要点：

\begin{itemize}
\item \textbf{spill 的代价主要不是流量}。一个 fp4 attention kernel 在
  D=512 时 spill 总量 234MB，但 local 流量仅 12.8GB/s——带宽上无关痛痒，真正的代价
  是\textbf{延迟暴露}：8 warps/SM 意味着每 scheduler 只有 2 个 warp 可切换，
  LDL 的几百 cycle 延迟藏不住，同时指令数也变多。
\item \textbf{寄存器预算可以算出来}。attention 的 O 累加器每线程占
  $k_{Br}\cdot D/256$ 个 f32 寄存器（M8N1 布局，$256=32$ lanes $\times 8$
  个 warp——consumer 全部 256 线程）；$k_{Br}{=}128$ 时为
  $D/2$，D=512 要 256 个——\textbf{撞穿 255 墙}。这正是大 D 必须换
  M4N2 atom（列拆给两个 warp，每线程 $D/4$）的数学来源（ch19）。
\end{itemize}

\subsection{Warp State：stall reasons，好 stall 与坏 stall}

warp 每个周期若不能发射，必有一个「停滞原因」。sm\_120 上的口径是
\texttt{smsp\_\allowbreak{}average\_\allowbreak{}warps\_\allowbreak{}issue\_\allowbreak{}stalled\_\allowbreak{}*\_\allowbreak{}per\_\allowbreak{}issue\_\allowbreak{}active.\allowbreak{}ratio}
——每个发射周期停在该原因上的平均 warp 数。表~\ref{tab:00-stall}
是实战中最常读的八项：

\begin{table}[htbp]
\centering
\caption{warp stall 指标字典（sm\_120 实测口径）}
\label{tab:00-stall}
\begin{tabular}{lp{5.2cm}p{5.6cm}}
\toprule
指标后缀 & 含义 & 高值时怎么办 \\
\midrule
\texttt{wait} & 固定延迟依赖（寄存器/ALU 结果未就绪） & 依赖链问题：拆链、增并行 warp、重排指令；attention 的 softmax$\to$PV 串行链是典型来源 \\
\texttt{long\_\allowbreak{}scoreboard} & 等 gmem/TMA 数据 & 装载不足：\texttt{cp.async}/TMA、加深 stage、prefetch \\
\texttt{short\_\allowbreak{}scoreboard} & 等 smem（MIO 队列） & LDS 链太长：swizzle 布局、减少 smem 往返 \\
\texttt{math\_\allowbreak{}pipe\_\allowbreak{}throttle} & 数学管线排队满 & \textbf{好信号}：计算饱和。保持或进一步减 FLOPs \\
\texttt{sleeping} & warp 休眠（mbarrier 退避） & \textbf{纯浪费}：producer/consumer 失衡或空转自旋，查 WS 结构 \\
\texttt{barrier} & 等显式栅栏 & 同步过频：减少 \texttt{\_\allowbreak{}\_\allowbreak{}syncthreads} 点、缩 block \\
\texttt{not\_\allowbreak{}selected} & 就绪但调度器选了别人 & \textbf{健康}：有冗余并行度可隐藏延迟 \\
\texttt{selected} & 正在发射（非停滞） & 基线，越大越好 \\
\bottomrule
\end{tabular}
\end{table}

读 stall 剖析结果的关键是\textbf{分好坏}：\texttt{math\_\allowbreak{}pipe\_\allowbreak{}throttle} 与
\texttt{not\_\allowbreak{}selected} 高说明机器在干活；\texttt{sleeping} 高几乎总是结构
问题；\texttt{wait} 与 \texttt{long\_\allowbreak{}scoreboard} 的\textbf{相对大小}直接
区分「依赖链受限」与「装载受限」——例见 \S\ref{sec:ch00-cases} 例三。

\subsection{Memory Workload：L1/L2/DRAM 分层归因}

三层各看两个数——命中率与吞吐：

\begin{itemize}
\item \textbf{L1/TEX}：\texttt{l1tex\_\allowbreak{}\_\allowbreak{}t\_\allowbreak{}sector\_\allowbreak{}hit\_\allowbreak{}rate.pct}。低于
  20\% 且数据本该复用时，考虑 tiling 进 smem 或改访问模式。注意 smem 与 L1
  \textbf{共池}（GB202 实测每多占 16KB smem，kernel 约慢 4--5$\mu$s），
  「多分 smem」不是免费的。
\item \textbf{L2}：\texttt{lts\_\allowbreak{}\_\allowbreak{}t\_\allowbreak{}sector\_\allowbreak{}hit\_\allowbreak{}rate.\allowbreak{}pct}。大 GEMM 上低
  命中可用 CTA swizzle / StreamK 提数据局部性；推理场景 K/V 重复读高命中
  是常态。
\item \textbf{DRAM}：\texttt{dram\_\allowbreak{}\_\allowbreak{}throughput}（峰值百分比）。贴满
  （$>$70\%）说明带宽墙已到，优化方向只剩减流量（混合精度、融合、去
  物化）。一条量化前处理链实测贴到饱和带宽 1.04TB/s（约为理论峰值
  1344GB/s 的 77\%）——此时「融合省 launch」已无杠杆，只有减流量才有，这个
  判断直接由 DRAM 指标给出。
\item \textbf{load/store 与合并访问}：global load 的每请求扇区数
  （\texttt{l1tex} 前缀的
  \texttt{average\_\allowbreak{}t\_\allowbreak{}sectors\_\allowbreak{}\_\allowbreak{}per\_\allowbreak{}request\_\allowbreak{}\allowbreak{}...} 系列指标）
  衡量合并访问质量——32-bit 连续访问的理想值是 4（一个 warp 的 128B 恰好
  4 个 32B sector）；明显超标说明访存发散，流量被放大，对策是内维连续、
  向量化（\texttt{float4} 128-bit 访存，ch03）与 AoS$\to$SoA。IO 压力的
  完整画像由三个正交维度拼出：\textbf{吞吐}（哪层饱和）、\textbf{命中率}
  （数据该由哪层负责）、\textbf{合并度}（流量有没有被访问模式放大）。
\end{itemize}

分层归因顺序：DRAM 饱和先看是不是 L2 miss 推的（L2 hit 低 + DRAM 高 =
工作集溢出），L2 高命中但 L1 低命中 = smem/L1 复用不足。IO 压力的「层」
定位错了，对策就全错。

\subsection{bank conflict 与 smem 指标}

\texttt{l1tex\_\allowbreak{}\_\allowbreak{}data\_\allowbreak{}bank\_\allowbreak{}conflicts\_\allowbreak{}pipe\_\allowbreak{}lsu\_\allowbreak{}mem\_\allowbreak{}\_\allowbreak{}shared\_\allowbreak{}op\_\allowbreak{}ld/\allowbreak{}st.sum}
直接数冲突次数（ch12 表 12-1 的数据来源）。判读要点：\textbf{先看占比再
动手}——曾有一个「消除 O 写回 bank conflict」的优化立项，NCU 一查
ld conflict 仅 0.39\%，且 MUFU 等待被 MMA 掩盖，整个方向当场否决。
冲突计数要换算成「每 warp 每 cycle 的额外重放」才有体感，绝对数大不等于
值得修。

\subsection{指令 mix 与 tensor pipe 利用率}

各执行管线的活跃占比（\texttt{pipe\_fma} / \texttt{pipe\_alu} /
\texttt{pipe\_lsu} / \texttt{pipe\_tensor}）回答「指令都花在哪了」。
对本书 Part II--IV 的主线而言，\textbf{tensor pipe 利用率就是 Tensor Core
饱和度}：fp16 GEMM 类 kernel 若 \texttt{pipe\_fma} 高而 \texttt{pipe\_tensor}
低，说明还在用 CUDA Core 算乘加，应该换 MMA 指令（ch10--11 的主题）。
一个 fp8 kernel 的实例：tensor pipe 从 49\% 提到 59.5\% 的那次优化，
判据正是「\texttt{math\_\allowbreak{}pipe\_\allowbreak{}throttle} 降了、tensor pipe 升了」——两条
指标共同确认依赖链被拆开、Tensor Core 被喂得更满。

\section{瓶颈判断决策树}\label{sec:ch00-tree}

把上述指标串成一条判断流程（数值阈值为工程经验值，非硬边界）：

\begin{figure}[H]
\centering
\begin{tikzpicture}[
  font=\footnotesize,
  box/.style={draw, rounded corners=2pt, align=center, inner sep=3pt,
    text width=2.6cm, minimum height=8mm},
  dec/.style={box, draw=tkBlue, fill=tkFillBlue},
  memb/.style={box, draw=tkOrange, fill=tkFillOrange},
  cpub/.style={box, draw=tkPurple, fill=tkFillPurple},
  latb/.style={box, draw=tkRed, fill=tkFillRed},
  balb/.style={box, draw=tkGreen, fill=tkFillGreen},
  det/.style={draw, align=left, inner sep=3pt, font=\scriptsize,
    text width=5.1cm, fill=white},
  detm/.style={det, draw=tkOrange!85},
  detc/.style={det, draw=tkPurple!85},
  detl/.style={det, draw=tkRed!70},
  badge/.style={circle, draw=black!70, fill=white, inner sep=0pt,
    minimum size=3.6mm, font=\scriptsize\bfseries, anchor=south east},
  ar/.style={-{Stealth[length=1.8mm]}, semithick},
  arm/.style={ar, draw=tkOrange},
  arc/.style={ar, draw=tkPurple},
  arl/.style={ar, draw=tkRed!85}
]
  \node[dec] (root) at (0, 0.55) {读 SOL\\（SM\% / MEM\%）};
  \node[badge] at (root.north west) {1};
  \node[memb] (mcond) at (3.8, 4.2) {MEM\% $\gg$ SM\%\\（差 $>$20）};
  \node[memb] (mver) at (7.9, 4.2) {memory-bound\\分层归因};
  \node[badge] at (mver.north west) {4};
  \node[detm] (md1) at (12.9, 5.4) {DRAM $>$70\% $\to$ 带宽墙\\对策：减流量（融合 / 低精度 / 去物化）};
  \node[detm] (md2) at (12.9, 4.2) {L2 hit $<$50\% 且 DRAM $>$40\%\\$\to$ 工作集 / L2 局部性（CTA swizzle）};
  \node[detm] (md3) at (12.9, 3.0) {L1 hit $<$20\% 而 L2 高\\$\to$ smem tiling / 复用不足};
  \node[cpub] (ccond) at (3.8, 1.2) {SM\% $\gg$ MEM\%};
  \node[cpub] (cver) at (7.9, 1.2) {compute-bound\\查饱和质量};
  \node[badge] at (cver.north west) {5};
  \node[detc] (cd1) at (12.9, 1.8) {tensor pipe 高 = 健康\\$\to$ 减 FLOPs 或更低精度};
  \node[detc] (cd2) at (12.9, 0.6) {pipe\_fma / alu 高、tensor 低\\$\to$ 该上 MMA 指令（ch10--11）};
  \node[latb] (lcond) at (3.8, -1.8) {双低\\（都 $<$40）};
  \node[latb] (lver) at (7.9, -1.8) {latency-bound\\stall 前二名分流};
  \node[badge] at (lver.north west) {2};
  \node[detl] (ld1) at (12.9, -0.6) {\texttt{wait} / \texttt{math\_pipe\_throttle} 主导\\$\to$ 拆链、增并行 warp};
  \node[detl] (ld2) at (12.9, -1.8) {\texttt{long\_scoreboard} 主导\\$\to$ 异步装载（TMA / cp.async）};
  \node[detl] (ld3) at (12.9, -3.0) {\texttt{sleeping} / \texttt{barrier} 主导\\$\to$ 重排 producer/consumer（WS）};
  \node[detl, dashed, text width=5.5cm] (occ) at (7.9, -4.55) {先查 occupancy 限制源：\\limiter = regs/smem $\to$ 先解资源\\（缩 tile、降 stage、launch\_bounds）\\「增并行 warp」以有 warp 可增为前提};
  \node[badge] at (occ.north west) {3};
  \node[balb] (dcond) at (3.8, -6.15) {双高\\（都 $>$60）};
  \node[balb] (dver) at (7.9, -6.15) {接近平衡上限\\$\to$ 停止优化};
  \draw[ar, draw=tkBlue] (root.east) -- (mcond.west);
  \draw[ar, draw=tkBlue] (root.east) -- (ccond.west);
  \draw[ar, draw=tkBlue] (root.east) -- (lcond.west);
  \draw[ar, draw=tkBlue] (root.east) -- (dcond.west);
  \draw[arm] (mcond) -- (mver);
  \draw[arc] (ccond) -- (cver);
  \draw[arl] (lcond) -- (lver);
  \draw[ar, draw=tkGreen] (dcond) -- (dver);
  \draw[arm] (mver.east) -- (md1.west);
  \draw[arm] (mver.east) -- (md2.west);
  \draw[arm] (mver.east) -- (md3.west);
  \draw[arc] (cver.east) -- (cd1.west);
  \draw[arc] (cver.east) -- (cd2.west);
  \draw[arl] (lver.east) -- (ld1.west);
  \draw[arl] (lver.east) -- (ld2.west);
  \draw[arl] (lver.east) -- (ld3.west);
  \draw[arl, dashed] (lver.south) -- (occ.north);
\end{tikzpicture}
\caption{瓶颈判断决策树：先读 SOL（SM\%/MEM\%）四分流——蓝为入口判断，橙为 memory-bound 链，紫为 compute-bound 链，红为 latency-bound 链，绿为双高（停止优化）；对策框白底、边框沿用所属分支色。节点角标 1--5 对应正文 enumerate 步骤（阈值为工程经验值）。}
\label{fig:0-1}
\end{figure}

\begin{enumerate}
\item \textbf{先看 SOL}：MEM\% $\gg$ SM\%（差 $>$20）$\to$ memory-bound，
  进第 4 步；SM\% $\gg$ MEM\% $\to$ compute-bound，进第 5 步；\textbf{双低}
  （都 $<$40）$\to$ latency/依赖链问题，进第 2 步；双高（都 $>$60）$\to$
  已接近平衡上限，停止优化。
\item \textbf{latency-bound 分流}（看 stall 前二名）：\texttt{wait} 或
  \texttt{math\_\allowbreak{}pipe\_\allowbreak{}throttle} 主导 $\to$ 依赖链/发射受限，对策是拆链
  （拆 CTA、拆 warp 分工）或增并行 warp；\texttt{long\_\allowbreak{}scoreboard} 主导
  $\to$ 数据供给不足，对策是异步装载（\texttt{cp.async}/TMA）与加深流水；
  \texttt{sleeping}/\texttt{barrier} 主导 $\to$ 同步结构问题，对策是重排
  producer/consumer 与栅栏。
\item \textbf{查 occupancy 是否是限制源}：若 achieved 低且 limiter 是
  registers/smem，先解决资源（缩 tile、降 stage、\texttt{\_\allowbreak{}\_\allowbreak{}launch\_\allowbreak{}bounds\_\allowbreak{}}），
  因为第 2 步的「增并行 warp」以有 warp 可增为前提。
\item \textbf{memory-bound 分层}：DRAM $>$70\% $\to$ 带宽墙，减流量；
  L2 hit $<$50\% 且 DRAM $>$40\% $\to$ 工作集/L2 局部性问题；L1 hit $<$20\%
  而 L2 高 $\to$ smem tiling/复用不足。顺带查 coalescing 与 bank conflict。
\item \textbf{compute-bound 查饱和质量}：tensor pipe 高 = 健康，优化转向
  减 FLOPs 或换更低精度；tensor pipe 低而 fma/alu 高 = 该上 MMA；指令数
  本身异常大时查发散与编译器胶水指令。
\end{enumerate}

\begin{table}[htbp]
\centering
\caption{症状 $\to$ 证据 $\to$ 结论 $\to$ 对策速查}
\label{tab:00-tree}
\begin{tabular}{p{3.4cm}p{4.2cm}p{2.6cm}p{4.4cm}}
\toprule
症状 & 指标证据 & 结论 & 对策 \\
\midrule
e2e 慢但不知谁慢 & nsys kern\_sum 占比 & 定位目标 kernel & 只对 Top-1/2 跑 ncu \\
kernel 慢、带宽没吃满 & MEM\%$\le$40，wait/mpt 主导 & 依赖链受限 & 拆链、2 CTA/SM、增并行 warp \\
装载等数据 & long\_scoreboard 高 & 供给不足 & TMA/cp.async、加深 stage \\
并行度上不去 & occupancy 低 + limiter=regs & 寄存器墙 & 缩 tile、M4N2 拆列、launch\_bounds \\
local ld/st 非零 & LDL/STL、local bytes & spill & 降每线程状态、查 255 墙 \\
空转严重 & sleeping 高 & WS 失衡 & 重排 producer/consumer \\
带宽贴顶 & DRAM $>$70\% & 带宽墙 & 减流量：融合/低精度/去物化 \\
smem 访问慢 & bank conflicts 高 & 冲突重放 & swizzle/pad（ch12） \\
TC 没用上 & pipe\_tensor 低 & CUDA Core 算矩阵 & 换 mma/wgmma 指令 \\
\bottomrule
\end{tabular}
\end{table}

\section{实战三例：占比、寄存器与 stall}\label{sec:ch00-cases}

\subsection{例一：nsys 先行——占比决定优化预算}

场景：一个 fp8/fp4 量化 attention 库的 e2e（含量化前处理链）对标
业界实现。e2e 一度落后，最初的直觉是「辅助链 kernel 太多（每 call
约 13 个）应该融合成一个 Mega Kernel」。先用全局视角（nsys + 每 kernel
带宽核算）确认：大 N 场景\textbf{主 kernel 占 e2e 时间 86--92\%}，而 aux
链每个 kernel 都已贴实测饱和带宽 1.04TB/s——融合只能省 launch 不能减流量，
大 N wall 差实测 $<$0.2\%，方向证伪；真正的杠杆在主 kernel 与引擎层
overlap。\textbf{方法论}：占比数据先给「优化预算上限」（14\% 的池子里抠
不出 30\% 的收益），带宽核算再判「这个池子还抠不抠得动」。

\subsection{例二：寄存器压力——255 墙、spill 与 setmaxnreg 的取舍}

场景：fp8 persist-D kernel 采用 warp specialization（WS）：1 个 producer
warp group（128T，专职 TMA）+ 1 个 consumer（256T，全计算），用
\texttt{setmaxnreg} 把寄存器从 producer（32 个）再分配给 consumer（232 个）。
寄存器压力的证据链分三层：

\begin{itemize}
\item \textbf{静态}：\texttt{cuobjdump -{}-dump-resource-usage}（D=768，
  sm\_120f）——fp8 split\_d REG:255 + STACK 1104--1944B（重 spill）；fp8
  m4n2 REG:255 + STACK 0--264B（几乎无 spill）；fp16 split\_d REG:196--254
  + STACK 约 1536B。STACK 字节就是溢出到 local 的体积。
\item \textbf{动态}：NCU 确认 \texttt{Block Limit Registers=1}（REG:255 的
  直接后果），理论 occupancy 25\%（12 warps/48）且 smem 守卫恒空转——限制源唯一。
\item \textbf{取舍}：consumer 想要更多寄存器（\texttt{setmaxnreg} 上限
  232，实测 224--232 区间性能平坦），但 O 累加器 $D/2$（M8N1）在 D=512 时
  要 256 个——\textbf{撞穿 255 墙}，于是 split-D 的 WS 变体被整体禁用，
  大 D 改走 M4N2 把每线程压力降到 $D/4$。这是「寄存器文件利用率 vs
  occupancy vs 算法结构」三方权衡的标准案例：不是把寄存器降下来就赢，
  而是选一个让 tile、并行度与 spill 同时可接受的工作点。
\end{itemize}

（注意 \texttt{setmaxnreg} 在 sm\_120a 上会被 ptxas C7506 静默忽略，
必须用 sm\_120f 构建，见 ch14。）

\subsection{例三：stall 证据链——好 stall 与坏 stall 的判读}

场景：sm\_89 fp8 persist-D 909.4$\mu$s vs 业界实现 710$\mu$s
（慢 28\%）。NCU 剖析结果：

\begin{itemize}
\item \texttt{math\_\allowbreak{}pipe\_\allowbreak{}throttle} \textbf{1.32}（第一名）、
  \texttt{wait} 0.82、\texttt{long\_\allowbreak{}scoreboard} \textbf{0.04}、tensor
  pipe 49\%、255 regs、1 CTA/SM。
\item 判读：\texttt{long\_\allowbreak{}scoreboard} 近零 $\Rightarrow$ 装载不是瓶颈，
  一切「加深流水/砍指令/缩 tile」无效（后续 6 项实验全部证伪，白跑的原因
  就是没先读这组数）；\texttt{mpt} 高但 tensor pipe 只有 49\% $\Rightarrow$
  管线在排队却没排满 Tensor Core——\textbf{依赖链/发射受限}。
\item 对策与验证：kBr=64 + kBc=64 + 128T（2 CTA/SM 拆链）$\to$ 864$\mu$s
  （$-5\%$），ncu 确认 2 blocks/SM、每线程寄存器降到 250；
  \textbf{指标闭环}：mpt 1.32$\to$0.72、tensor pipe 49\%$\to$59.5\%——
  「排队变短、TC 更饱」与 wall 时间下降互为印证。
\end{itemize}

同族的「坏 stall」案例：bias 注入的 smem tile 版本 \texttt{sleeping} 从
3.5\% 涨到 22.3\%（SASS 采样里 \texttt{SYNCS.\allowbreak{}PHASECHK} mbarrier 自旋占
18\%）——LDS 进了 consumer 关键路径、两级 QK 流水断供，全 warp 自旋等 K，
该版本被指标否决。

\section{常见坑}

\begin{enumerate}
\item \textbf{\texttt{-{}-kernel-name} 精确匹配陷阱}。模板 kernel 名带长
  后缀，前缀匹配报 \texttt{No kernels were profiled} 且生成空 report——
  一律 \texttt{regex:} 前缀。
\item \textbf{ncu harness 与 bench 配置不一致}。profiling 用的配置必须与
  bench 默认逐 flag 对齐；显式传 flag 会覆盖 bench 内部的默认值，实测逐
  flag 传参 vs 全默认可差 13\%，指标分析对象可能根本不是 bench 跑的那个
  kernel 形态。
\item \textbf{A/B 不校验公共参照}。GPU 级干扰让两轮\textbf{同比例膨胀}
  （约 $2\times$），对比出假加速。先核对 bench 输出里的 SDPA 参照时间
  （$\pm$1\% 内）再谈结论；抖动大的实现用 min-of-N。
\item \textbf{把 section 输出当 raw metric}。\texttt{-{}-csv} 默认输出是
  rule-based 汇总；要具体计数器数值必须 \texttt{-{}-page raw}（配合
  \texttt{grep -A3} 抓多行结构）。
\item \textbf{拿 ncu 的耗时报性能}。replay 多遍 + 锁频 + 串行化，报告内
  时间只做相对比较；性能数字一律回真实 bench 采。
\item \textbf{不标性能口径}。写结论必须标注 GPU 型号/架构/日期（PRO 5000
  sm\_120a 与 5090 sm\_120 的架构级结论可互证，绝对数字不可）；fallback
  也要查——小 seqlen（$<$512）时库可能自动 fallback 到 PyTorch SDPA，profile
  到的根本不是目标 kernel。
\item \textbf{nsys 全量采集撑爆磁盘}。大模型 e2e 无限制采集曾写爆
  \texttt{/tmp}（500GiB+）。TMPDIR 落项目目录、单域、最小 repeat、事后必清
  （\S\ref{sec:ch00-nsys}）。
\end{enumerate}

\section{小结}

本章没有写一行 kernel，却给出了全书其余 26 章共同的工作节奏：
\textbf{nsys 定位（谁慢）$\to$ ncu 定性（为什么慢）$\to$ 单变量修改 $\to$
复测 diff $\to$ 指标闭环}。指标字典里最值得带走的三句话：SOL 双低不是
「没病」而是 latency-bound，答案在 stall 剖析结果里；stall 分好坏，
\texttt{math\_\allowbreak{}pipe\_\allowbreak{}throttle} 高是 Tensor Core 喂满的勋章，\texttt{sleeping}
高是结构问题的警报；寄存器 255 墙、spill 与 occupancy 是同一份预算的三种
花法，读数时要一起读。工具路径：ncu/nsys 随 CUDA Toolkit 安装（通常在
\texttt{/\allowbreak{}usr/\allowbreak{}local/\allowbreak{}cuda/\allowbreak{}bin}），不在 PATH 时用绝对路径调用。

\section*{延伸阅读}
\begin{itemize}
\item NVIDIA Nsight Systems 文档：
  \href{https://developer.nvidia.com/nsight-systems}{developer.nvidia.com/nsight-systems}
\item NVIDIA Nsight Compute 文档（含 Profiling Guide 的完整指标索引）：
  \href{https://developer.nvidia.com/nsight-compute}{developer.nvidia.com/nsight-compute}
\item 本书各章「性能分析」节：ch09--14（GEMM/TMA/WS 的指标判读）、ch12
  （bank conflict 计数器实战）、ch19（大 D 寄存器压力模型）、ch01
  （roofline 与 occupancy 的理论侧）。
\end{itemize}
